\section{The limit of the discriminant function and the inconsistency}
\label{sec:disc}

\subsection{The limit of the discriminant function}

Let $x_0\in\pi_1$ and consider the rescaled new observation
\begin{equation}
u_0=\frac{x_0-(\mu_1+\mu_2)/2}{\sqrt{\theta}}
=\eps_0\frac{\mu}{2\sqrt{\theta}}+\frac{v_0+\xi_0}{\sqrt{\theta}},
\quad\eps_0=+1,
\label{eq:6.1}
\end{equation}
where $v_0=\sum_{s\le m_1}\sqrt{\lambda_{1s}}z_{0s}h_{1s}$ and
$\xi_0=\sum_{s>m_1}\sqrt{\lambda_{1s}}z_{0s}h_{1s}$. From~\eqref{eq:5.4} we have
\begin{equation}
f(u_0)=\sum_{p\in\cI}\beta_p\gamma_p+b,
\quad\gamma_p=u_p^Tu_0.
\label{eq:6.2}
\end{equation}

\begin{theorem}
\label{thm:3}
Assume \textup{(A1)} to \textup{(A4)} and \textup{(S)}. Then, as $d\to\infty$,
it holds that
\begin{equation}
f(u_0)\wto f^*(z_0)
=\sum_{(i,j)\in\cI}\beta_{ij}^*\,g_{(ij),(0)}^*+b^*,
\label{eq:6.3}
\end{equation}
where $z_0=(z_{01},\dots,z_{0m_1})^T$, $(\beta^*,b^*)$ is given in
Theorem~\ref{thm:2}, and
\begin{equation}
\begin{aligned}
g_{(ij),(0)}^*
&=\frac{\eps_i\delta}{4}
+\frac{\eps_i\sqrt{\delta}}{2}\sum_{t\le m_1}\sqrt{c_{1t}}\,\beta_{1t}z_{0t}
+\frac{\sqrt{\delta}}{2}\sum_{s\le m_i}\sqrt{c_{is}}\,\beta_{is}z_{ijs}\\
&\qquad+\sum_{s\le m_i}\sum_{t\le m_1}\sqrt{c_{is}c_{1t}}\,r_{st}^{(i1)}z_{ijs}z_{0t}.
\end{aligned}
\label{eq:6.4}
\end{equation}
\end{theorem}

We emphasize that the term $c_{10}$ does not appear in~\eqref{eq:6.4}, that is, the
non-spiked part of $x_0$ does not contribute to the limit.

\begin{proof}
We follow the same steps as in the proof of Theorem~\ref{thm:1}, replacing
$u_{kl}$ by $u_0$. The only difference is that $x_0$ is independent of all the
training observations, and this single difference changes the conclusion.

\textit{Step~1.} From~\eqref{eq:4.1} and~\eqref{eq:6.1}, the bilinearity gives
\begin{equation}
\gamma_p=u_{ij}^Tu_0
=\underbrace{\frac{\eps_i\eps_0\Delta}{4\theta}}_{\mathrm{(I)}}
+\underbrace{\frac{\eps_i}{2}\frac{\mu^T(x_0-\mu_1)}{\theta}}_{\mathrm{(II)}}
+\underbrace{\frac{\eps_0}{2}\frac{\mu^T(x_{ij}-\mu_i)}{\theta}}_{\mathrm{(III)}}
+\underbrace{\frac{(x_{ij}-\mu_i)^T(x_0-\mu_1)}{\theta}}_{\mathrm{(IV)}},
\label{eq:6.5}
\end{equation}
where $p=(i,j)$ and we used $\eps_0=+1$ in~(I).

\textit{Step~2.} The term~(I) does not involve any random variable, and from~(A3)
we have $\eps_i\delta/4$, which is the first term of~\eqref{eq:6.4}.

\textit{Step~3.} Since $x_0\in\pi_1$, we may apply Lemma~\ref{lem:3} with $i=1$ to
$x_0$. Note that Lemma~\ref{lem:3} is a statement about a single observation and
does not use any relation to the other observations. Hence
\[
\frac{\mu^T(x_0-\mu_1)}{\theta}
\longrightarrow\sqrt{\delta}\sum_{t\le m_1}\sqrt{c_{1t}}\,\beta_{1t}z_{0t}.
\]
Multiplying by $\eps_i/2$, we obtain the second term of~\eqref{eq:6.4}.

\textit{Step~4.} Applying Lemma~\ref{lem:3} to $x_{ij}$ and using $\eps_0=+1$, we
obtain the third term of~\eqref{eq:6.4}.

\textit{Step~5.} We consider the term~(IV), which is the essential difference from
Theorem~\ref{thm:1}. Since $x_0$ is independent of all the training observations
$\{x_{ij}\}$, the term~(IV) is always the inner product of two distinct independent
observations, wherever $(i,j)$ is taken in $\cI$. The case of the inner product of
an observation with itself, which appeared in the diagonal components in
Theorem~\ref{thm:1}, never occurs here. We divide into two cases.

(IV-a) The case $i=1$. Although $x_{1j}$ and $x_0$ are observations from the same
population $\pi_1$, they are independent, so that we may apply~\eqref{eq:3.1} of
Lemma~\ref{lem:1}. We note that the condition $j\neq k$ in Lemma~\ref{lem:1} was
used only to guarantee the independence. Hence
\[
\frac{(x_{1j}-\mu_1)^T(x_0-\mu_1)}{\theta}
=\sum_{s=1}^{m_1}\frac{\lambda_{1s}}{\theta}z_{1js}z_{0s}+o_P(1).
\]
This agrees with the fourth term of~\eqref{eq:6.4} with $i=1$, because
$r_{st}^{(11)}=\delta_{st}$.

(IV-b) The case $i=2$. Since $x_{2j}$ and $x_0$ are independent observations from
different populations, we may apply Lemma~\ref{lem:2}. Exchanging the roles of the
two factors, we obtain
\[
\frac{(x_{2j}-\mu_2)^T(x_0-\mu_1)}{\theta}
=\sum_{s=1}^{m_2}\sum_{t=1}^{m_1}\frac{\sqrt{\lambda_{2s}\lambda_{1t}}}{\theta}
\rho_{ts}\,z_{2js}z_{0t}+o_P(1),
\]
where $\rho_{ts}=h_{2s}^Th_{1t}$. This agrees with the fourth term of
\eqref{eq:6.4} with $i=2$ under the convention $r_{st}^{(21)}=r_{ts}$.

\textit{Step~6.} We explain why $c_{10}$ does not appear. In Theorem~\ref{thm:1},
the term $c_{i0}$ came from~\eqref{eq:3.2}, which was used when $(i,j)=(k,l)$. The
corresponding term here is $\xi_{ij}^T\xi_0/\theta$. Since $\xi_{ij}$ and $\xi_0$
are independent, we have $E(\xi_{ij}^T\xi_0)=0$ and, by the same computation as in
\eqref{eq:3.3},
\[
E[(\xi_{ij}^T\xi_0)^2]=\tr(\Sigma_{i*}\Sigma_{1*})=o(\theta^2),
\]
so that $\xi_{ij}^T\xi_0/\theta=o_P(1)$ from Chebyshev's inequality.

Geometrically, the noise direction $\xi_0$ of the new observation is asymptotically
orthogonal to the noise directions $\xi_{ij}$ of the training data, because two
independent random directions are nearly orthogonal in high dimensions. On the
other hand, $\xi_{ij}$ is of course not orthogonal to itself, so that $c_{i0}$
remains only in the diagonal components in Theorem~\ref{thm:1}. Since the
discriminant function uses only the cross terms $u_p^Tu_0$, the noise of $x_0$
does not affect the classification.

\textit{Step~7.} From Steps~2 to~5, for each $p\in\cI$, the quantity $\gamma_p$ is
written as a continuous function of $Z^{(d)}$ plus $o_P(1)$. Since the number of
the components is $N$, which does not depend on $d$, we obtain
\begin{equation}
\gamma=(\gamma_p)_{p\in\cI}\wto\gamma^*=(g_{(ij),(0)}^*)_{(i,j)\in\cI}
\label{eq:6.6}
\end{equation}
from~(A4) together with the continuous mapping theorem and Slutsky's theorem.

\textit{Step~8.} We show the joint convergence of $(G,\gamma)$. The $N^2$
components of $G$ given in Theorem~\ref{thm:1} and the $N$ components of $\gamma$
given in Step~7 are all written as continuous functions of the common random
vector $Z^{(d)}$ plus $o_P(1)$. That is, there exists a continuous map $\Psi$,
which depends on the coefficients $c_{is}$, $\beta_{is}$ and $r_{st}$ converging
deterministically, such that $(G,\gamma)=\Psi(Z^{(d)})+o_P(1)$ uniformly on
compact sets. From~(A4) we have $Z^{(d)}\wto Z$, so that
\begin{equation}
(G,\gamma)\wto(G^*,\gamma^*)=\Psi(Z)
\label{eq:6.7}
\end{equation}
jointly. We note that this step is necessary, because the marginal convergences
shown separately do not imply the joint convergence.

\textit{Step~9.} We consider the map
\[
\Xi(G,\gamma)=\sum_p\beta_p(G)\gamma_p+b(G).
\]
From Lemma~\ref{lem:5}, the map $G\mapsto(\beta(G),b(G))$ is continuous on
$\cG_+$, and the inner product is continuous, so that $\Xi$ is continuous on
$\cG_+\times\R^N$. From~\eqref{eq:6.2} we have $f(u_0)=\Xi(G,\gamma)$. From
Theorem~\ref{thm:1} we have $P(G^*\succ O)=1$, so that the set of the discontinuity
points of $\Xi$ has measure zero under the limiting distribution. Thus we obtain
\eqref{eq:6.3} from~\eqref{eq:6.7} and the continuous mapping theorem.
\end{proof}

\begin{table}[ht]
\centering
\caption{Correspondence between the proofs of Theorems~\ref{thm:1} and~\ref{thm:3}.}
\label{tab:1}
\begin{tabular}{@{}lll@{}}
\toprule
Term & Theorem~\ref{thm:1} (training vs training)
& Theorem~\ref{thm:3} (training vs new) \\
\midrule
(I) deterministic
& $\eps_i\eps_k\delta/4$
& $\eps_i\delta/4$ (since $\eps_0=1$) \\
(II), (III) linear
& Lemma~\ref{lem:3} on both sides
& Lemma~\ref{lem:3} on both sides, one being $z_0$ \\
(IV), distinct
& \eqref{eq:3.1} or Lemma~\ref{lem:2}
& always this case \\
(IV), identical
& \eqref{eq:3.2}, giving $+c_{i0}$
& never occurs, no $c_{i0}$ \\
\bottomrule
\end{tabular}
\end{table}

\subsection{The inconsistency}

We now derive the inconsistency of the SVM. Here we need to derive a lower bound of
the expectation from the convergence in distribution. For this purpose we use
Fatou's lemma, which we summarize in Appendix~\ref{app:A.2} together with its
intuitive meaning. Roughly speaking, the lemma states that a non-negative mass may
escape in the limiting procedure but never emerges, and this is exactly the
direction of the inequality which we need.

\begin{corollary}
\label{cor:1}
Assume \textup{(A1)} to \textup{(A4)} and \textup{(S)}. Assume that $z_0$ has a
continuous distribution whose support is the whole space. Let
$e(1)=P\{f(u_0)<0\mid\text{training data}\}$. Then it holds that
\[
e(1)\wto e^*(1)=P_{z_0}\{f^*(z_0)<0\}.
\]
Moreover, if $f^*$ is a non-constant function of $z_0$ with positive probability,
then
\[
\liminf_{d\to\infty}E\{e(1)+e(2)\}>0,
\]
that is, the hard-margin SVM and the BC-SVM do not hold the consistency property.
\end{corollary}

\begin{proof}
\textit{Step~1.} From Theorem~\ref{thm:3} we have $f(u_0)\wto f^*(z_0)$. Since
$z_0$ is independent of the training data and the distribution of $f^*(z_0)$ is
continuous, the boundary $\{f^*=0\}$ has measure zero under the limiting
distribution. Hence $e(1)\wto e^*(1)$ from the Portmanteau theorem.

\textit{Step~2.} From~\eqref{eq:6.3} and~\eqref{eq:6.4}, $f^*(z_0)$ is a polynomial
in $z_0$ consisting of an affine part and a quadratic part. If it is non-constant,
then it takes negative values on a set of positive Lebesgue measure, and hence
$E\{e^*(1)\}>0$ because the support of $z_0$ is the whole space.

\textit{Step~3.} Since $e(1)\ge 0$, we may apply Lemma~\ref{lem:7} in
Appendix~\ref{app:A.2} and obtain
\[
E\{e^*(1)\}\le\liminf_{d\to\infty}E\{e(1)\}.
\]
Adding $e(2)$ does not change the inequality.
\end{proof}

\begin{remark}
\label{rem:5}
We emphasize that the inconsistency in Corollary~\ref{cor:1} is different from the
strong inconsistency of \citet{nakayama2017}. They showed that the
misclassification rate tends to one when the sample sizes are imbalanced, which is
caused by the deterministic bias $\kappa$ in~\eqref{eq:1.3}. Here the
misclassification rate tends neither to zero nor to one, and it is caused by the
randomness which survives in the limit. Hence the bias correction alone cannot
remove this difficulty. We discuss this point in Sections~\ref{sec:bias}
and~\ref{sec:scsvm}.
\end{remark}
